$latexcv-preamble.tex()$

%-------------------------------------------------------------------------------
%	CLASSIC
%
%	A centred title with a rule between the name and the document type, contact
%	details paired off beneath it, and centred section headings over full-width
%	entries. The plainest of the styles: no colour blocks, no sidebar.
%-------------------------------------------------------------------------------

% Upstream sets a negative bottom margin to pull its single page together. That
% is a one-page trick and misbehaves as soon as the CV runs on, so an ordinary
% bottom margin is used here.
$if(geometry)$
\geometry{$for(geometry)$$geometry$$sep$,$endfor$}
$else$
\geometry{top=1.25cm, bottom=1.5cm, left=1.5cm, right=1.5cm}
$endif$

%-------------------------------------------------------------------------------
%	HEADINGS
%-------------------------------------------------------------------------------
\newcommand{\cvsection}[1]{%
  \vspace{4pt}
  \begin{center}
    \large\textcolor{sectcol}{\textbf{#1}}
  \end{center}
  \vspace{-2pt}}

\newcommand{\cvsubsection}[1]{%
  \vspace{2pt}
  {\normalsize\textcolor{bgcol}{\textbf{#1}}}\par
  \vspace{2pt}}

% Upstream pairs off single lines of metadata under the title: a value on the
% left, a labelled value on the right. Kept for documents that want it in
% hand-written LaTeX.
\newcommand{\metasection}[2]{%
  \footnotesize{#2} \hspace*{\fill} \footnotesize{#1}\\[1pt]}

%-------------------------------------------------------------------------------
%	ENTRIES
%-------------------------------------------------------------------------------
% The headline shared by \vitaeentry and \vitaebrief: what/with/where set flush
% left, the date flush right in the accent colour. tabular* with a computed
% first column keeps the date hard against the right margin whatever the page
% geometry is, rather than at upstream's fixed 13.6cm.
\newcommand{\cventryhead}[4]{%
  \begin{tabular*}{\linewidth}{@{}p{\dimexpr\linewidth-\cvwhenwidth-2\tabcolsep\relax}@{\extracolsep{\fill}}x{\cvwhenwidth}@{}}
    \textbf{#2}%
    \ifthenelse{\isempty{#3}}{}{ - \textcolor{bgcol}{#3}}%
    \ifthenelse{\isempty{#4}}{}{, \textit{\textcolor{bgcol}{#4}}}%
    & \vspace{2.5pt}\textcolor{sectcol}{#1}
  \end{tabular*}}

\newcommand{\vitaeentry}[5]{%
  \cventryhead{#1}{#2}{#3}{#4}
  \vspace{-8pt}
  \textcolor{softcol}{\hrule}
  \vspace{6pt}
  #5
  \vspace{3pt}}

% No rule and no elaboration: a brief entry is one line, and a hairline under
% every line would read as a table.
\newcommand{\vitaebrief}[4]{%
  \cventryhead{#1}{#2}{#3}{#4}
  \vspace{2pt}}

% The measurement and the row itself are in the shared preamble: every design
% sets this listing the same way, and only the floor below differs.
\newcommand{\vitaeskill}[2]{\cvskillrow{#1}{#2}}

% Upstream marks its elaboration with a centred dot. The list form is used here
% so that a long point wraps with a hanging indent instead of back to the
% margin. \textperiodcentered rather than upstream's math-mode \cdot, since a
% literal dollar sign would have to be escaped from the Pandoc template too.
\newenvironment{vitaewhy}{%
  \begin{itemize}[label=\textperiodcentered, topsep=0pt, partopsep=0pt, itemsep=2pt,
                  parsep=0pt, leftmargin=1.2em, labelsep=0.5em]%
}{%
  \end{itemize}%
}

%-------------------------------------------------------------------------------
%	PROFILE PICTURE (optional)
%-------------------------------------------------------------------------------
% How wide the picture is set, needed as a length rather than as the template
% variable itself because the contact column beside it is the remainder of the
% measure and has to be computed from it.
\newlength{\cvpicwidth}
\newlength{\cvpicgap}
\setlength{\cvpicgap}{1.5em}

\begin{document}

$if(profilepic)$
\setlength{\cvpicwidth}{$if(profilepic-width)$$profilepic-width$$else$0.22\linewidth$endif$}
$endif$

%-------------------------------------------------------------------------------
%	TITLE
%-------------------------------------------------------------------------------
\vspace{-8pt}
\begin{center}
  \HUGE \textsc{$name$$if(surname)$ $surname$$endif$}%
  $if(docname)$\cvtitlerule\textsc{$docname$}$endif$\\[2pt]
  $if(position)$\small $position$$endif$
\end{center}

\vspace{6pt}

%-------------------------------------------------------------------------------
%	CONTACT DETAILS
%-------------------------------------------------------------------------------
$if(profilepic)$
% With a picture the head of the document reflows into two columns under the
% title: the photo on the left, the contacts stacked against the right margin
% beside it, one to a line. Centring the contacts under a picture that is itself
% off to one side would leave the header visibly lopsided, so the two are set as
% a pair that balances across the measure.
%
% Vertically centred on each other, so a document with three contacts and one
% with ten both sit level with the photo rather than hanging from its top edge.
\noindent
\begin{minipage}[c]{\cvpicwidth}
  \includegraphics[width=\cvpicwidth]{$profilepic$}
\end{minipage}%
\hfill
\begin{minipage}[c]{\dimexpr\linewidth-\cvpicwidth-\cvpicgap\relax}
  \raggedleft
  \footnotesize
  \renewcommand{\cvcontactsep}{\\}
  \cvcontacts
\end{minipage}

\vspace{6pt}
$else$
% Centred and wrapped rather than upstream's two hard columns, so that a
% document supplying three contacts and one supplying ten both come out
% balanced.
\begin{center}
  \footnotesize
  \renewcommand{\cvcontactsep}{\hspace{10pt}\textcolor{softcol}{|}\hspace{10pt}}
  \cvcontacts
\end{center}
$endif$

\vspace{-2pt}
\textcolor{softcol}{\hrule}
\vspace{6pt}

\normalsize

$latexcv-body.tex()$

%-------------------------------------------------------------------------------
%	FOOTER
%-------------------------------------------------------------------------------
% Upstream pins a footer to the bottom of its single page with \vfill. That
% would force a page break on a CV that already fills the page, so this is set
% as an ordinary trailing line and only when asked for.
$if(footer)$
\vspace{12pt}
\textcolor{softcol}{\hrule}
\vspace{4pt}
\begin{center}
  \footnotesize\textcolor{sectcol}{$footer$}
\end{center}
$endif$

\end{document}
